% Tipo de documento. En este caso es un art�culo, para folios A4, tama�o de la fuente 11pt y con p�gina separada para el t�tulo
\documentclass[a4paper,11pt,titlepage]{article}

% Carga de paquetes necesarios. OrdenesArticle es un paquete personalizado
\usepackage[spanish]{babel} 
\RequirePackage[T1]{fontenc}
\RequirePackage[ansinew]{inputenx} 
\usepackage[spanish,cap,cont,title,fancy]{OrdenesArticle}
\usepackage{lmodern}
\usepackage{array}
\usepackage{graphicx}
\usepackage{hyperref}
\usepackage{pifont}
\usepackage{listings}
\usepackage[usenames,dvipsnames]{color}
\usepackage{colortbl}
\usepackage{color}
\usepackage{ifthen}
\usepackage{longtable}
\usepackage{enumerate}
\hypersetup{bookmarksopen,bookmarksopenlevel=4,linktocpage,colorlinks,urlcolor=blue,citecolor=blue,
						linkcolor=blue,filecolor=blue,pdfnewwindow,
						pdftitle={PlanCoDE: Planificador Colaborativo para el Dise�o de Estrategias y Acciones de Emergencia},
						pdfauthor={Juan Andrada Romero, Jose Domingo L�pez L�pez},
						pdfsubject={Sistemas para la Colaboraci�n}}


% Macro para definir una lista personalizada 
\newenvironment{milista}%
{\begin{list}{\textbullet}%
{\settowidth{\labelwidth}{\textbullet} \setlength{\leftmargin}{\dimexpr\labelsep+\labelwidth+5pt}
\setlength{\itemsep}{\dimexpr 0.5ex plus 0.25ex minus 0.25ex}
\setlength{\parsep}{\itemsep}
\setlength{\partopsep}{\itemsep}
\addtolength{\topsep}{-7.5pt}
}}%
{\end{list}}

% Macro para insertar una imagen
%       Uso: \imagen{nombreFichero}{Factor escala}{Caption (leyenda)}{Label (identificador para referenciarla)}
% -------------------------------------------------------------------------------------------------------------
\def\imagen#1#2#3#4{
 \begin{figure}[h]
 \begin{center}
   \scalebox{#2}{\includegraphics{#1}}
 \caption {#3}
 \label{#4}
 \end{center}
 \end{figure}
 }
 

% Definicion del "listings" para el lenguaje Java
\lstdefinelanguage{Java}
{
 morecomment = [l]{//}, 
 morecomment = [l]{///},
 morecomment = [s]{/*}{*/},
 morestring = [b]", 
 sensitive = true,
 morekeywords = {package, static, while, switch, break, line, void, String, Object, int, Integer, instanceof, else, if, for, private, return, new, public, class, import, int, boolean, true, false, extends, final, super, protected, abstract, this, do, float, double, null, try, catch, implements}
}

% Configuraci�n del c�digo incluido para el lenguaje Java
\lstset{
  language=Java,
  basicstyle=\footnotesize,
  backgroundcolor=\color{white},
  showspaces=false,
  showstringspaces=false,
  showtabs=false,
  frame=single,
  tabsize=2,
  captionpos=b,
  breaklines=true,
  breakatwhitespace=false,
  escapeinside={\%},
  keywordstyle = \color [rgb]{0,0,1},
  commentstyle = \color [rgb]{0.133,0.545,0.133},
  stringstyle = \color [rgb]{0.627,0.126,0.941}
}

\begin{document}

% En las p�ginas de portada e �ndices, no hay encabezado ni pie de p�gina
\pagestyle{empty} 

% Se incluye la portada
\begin{titlepage}
	\begin{center}
  	{\LARGE UNIVERSIDAD DE CASTILLA-LA MANCHA} \\
  	\bigskip
  	{\Large ESCUELA SUPERIOR DE INFORM�TICA} \\
  	\vspace{20mm}
  	\includegraphics[scale=0.45, keepaspectratio]{./images/esi_bw.png} \\
  	\vspace{20mm}
  	{\Huge \textbf{Almacenamiento y Recuperaci�n de la Informaci�n}}\\
  	\vspace{10mm}
  	{\LARGE \textsc{\textbf{M�dulo de importaci�n e indexaci�n de documentos}}}\\
  	\vspace{30mm}
  	\today
  	\vspace{40mm}
  	\begin{flushleft}
  		{\large Juan Andrada Romero}\\
  		\vspace{1mm}
  		{\large Jose Domingo L�pez L�pez}\\		
  	\end{flushleft}
	\end{center}
\end{titlepage}
\clearpage

% Texto del reverso de la portada
%%\mbox{}
%%\vspace{18cm}
%%\begin{small}
% Se ajusta la separaci�n entre p�rrafos
%%\parskip=10pt 

%%\copyright~ 2008/2009 Juan Andrada Romero. Universidad de Castilla La Mancha, Escuela Superior de Inform�tica de Ciudad Real.

%%Se permite la modificaci�n, copia y distribuci�n de este documento, seg�n la licencia de documentaci�n GNU (\url{http://www.gnu.org}).

%%Este documento fue compuesto con \LaTeX{}. Im�genes generadas con Power Point y Gimp.
%%\end{small}

%%\newpage

% En las p�ginas de �ndices y prefacio, se utiliza numeraci�n romana
\pagenumbering{Roman}

% Se crea el �ndice
\tableofcontents
% Se pasa p�gina y se a�ade el �ndice de figuras al �ndice principal
%\clearpage\phantomsection
%\addcontentsline{toc}{section}{\listfigurename}
% Se crea el �ndice de figuras
%\listoffigures
% Si se quiere crear un �ndice de tablas se pondr�a: \listoftables

\newpage

% Se ajusta la separaci�n entre p�rrafos
\parskip=10pt

% Se a�ade el prefacio al �ndice
%\clearpage\phantomsection
%\addcontentsline{toc}{section}{Prefacio}

% Comienza el contenido del documento. Se utilizan n�meros ar�bigos y el encabezado y pie de p�gina personalizado
\pagenumbering{arabic}
\pagestyle{fancy}

\section{Introducci�n}


\section{�Qu� es PlanCoDE?}

\subsection{Diagrama de casos de uso}

\clearpage
\section{Clasificaci�n} \label{clasificacion}

En los siguientes apartados, se realizar� una clasificaci�n de la herramienta PlanCode, seg�n los diferentes autores que aparecen en \cite{groupware}.

\subsection{Dimensiones espacio-temporal}

Atendiendo a la clasificaci�n en tiempo y espacio propuesta por \textbf{Johansen}, PlanCode puede clasificarse en la categor�a \textit{diferente lugar, mismo tiempo}, es decir, es una herramienta de interacci�n s�ncrona distribuida (ver Figura \ref{fig:johansen}). Esto es as� porque los diferentes usuarios deben estar conectados en la misma sesi�n y al mismo tiempo para poder recibir los mensajes del chat, la imagen del mapa cargada por alguno de ellos y los trazos que se van realizando.

\imagen{./imagenes/johansen.png}{0.3}{Taxonom�a en tiempo y espacio (Johansen; 1991)}{fig:johansen}

Por otra parte, atendiendo a la clasificaci�n espacio-temporal propuesta por \textbf{Grudin}, la herramienta se puede clasificar tambi�n en la categor�a \textit{diferente lugar, mismo tiempo}, por las mismas razones que en el caso anterior. 


\subsection{Dominio de aplicaci�n}

Siguiendo la clasificaci�n del nivel de aplicaciones de \textbf{Ellis}, PlanCode cumple con las caracter�sticas de los siguientes sistemas:
\begin{enumerate}[A.]
	\item \textbf{Sistema de mensajes}: PlanCode se puede considerar un sistema de mensajes porque permite el intercambio de mensajes de diferente naturaleza entre los usuarios, como es el intercambio de mensajes de texto a trav�s del chat, im�genes cuando se carga un mapa y el intercambio de trazos libres.
\end{enumerate}

\subsection{Modelo 3C}

Si atendemos al objetivo principal de PlanCode (el dise�o de rutas sobre un mapa en situaciones de emergencias), dicha herramienta se puede clasificar como una herramienta de colaboraci�n, pues, como ya se ha comentado, los trazos pueden ser dibujados por cualquier usuario, apoyando al resto de trazos que ya est�n dibujados. Del mismo modo, cualquier usuario podr�a modificar la imagen del mapa que se est� mostrando si lo cree conveniente. \\
\indent Pero, como PlanCode cuenta con un chat, se puede considerar tambi�n una herramienta de comunicaci�n y coordinaci�n.



\clearpage
\section{Requisitos del sistema}

Para poder utilizar \textbf{FAQ Manager} ser� necesario disponer de un servidor de p�ginas Web (por ejemplo, Apache HTTP Server \cite{apache-homepage}) que pueda procesar c�digo PHP \cite{php-homepage}. Adem�s, ser� necesario un sistema gestor de bases de datos como MySQL \cite{mysql-homepage}. No obstante, para simplificar el proceso de instaci�n y configuraci�n de todo este software, se recomienda el uso de WAMP \cite{wamp-homepage}.

Una vez instalado y configurado el software necesario, deber� editar el fichero \texttt{config.php} que se incluye en la distribuci�n de \textit{FAQ Manager}. En este fichero, de estructura similiar a la de Figura \ref{fig:config}, deber� cumplimentar debidamente el dominio sobre el que se aloja el sistema gestor de bases de datos, el nombre de la base de datos, el nombre de un usuario de la base de datos con los permisos b�sicos para realizar operaciones \textit{CRUD} (Create, Read, Update, Delete) y su contrase�a.

\imagen{./imagenes/config.png}{0.7}{Fichero \texttt{config.php}}{fig:config}

\clearpage
\input{5dise�o.tex}
\clearpage
\section{Manual de usuario} \label{manual}

\subsection{Requisitos del sistema}

El sistema PlanCode se ha desarrollado bajo el lenguaje de programaci�n Java para hacer posible su portabilidad entre distintos sistemas, siempre y cuando �stos tengan instalada una m�quina virtual de Java. De este modo, las aplicaciones son compatibles tanto en entornos basados en Windows 98/NT/XP/Vista/Seven como en entornos GNU/Linux.

\subsubsection{Entorno de Java}
Para comprobar si ya dispone de una instalaci�n del entorno de ejecucic�n de Java en el equipo y que est� configurado correctamente, abra una ventana de comandos \textit{(Inicio -\textgreater Ejecutar -\textgreater cmd)} y escriba ''java -version''. Debe obtener un resultado similar al de la Figura \ref{fig:java-version}.

\imagen{./imagenes/usermanual/java-version.png}{0.80}{C�mo averiguar la versi�n de Java}{fig:java-version}

Si aparece un mensaje de error, pruebe lo siguiente: 
\begin{milista}
	\item Compruebe si tiene instalado Java. Para ello compruebe que exista el directorio Java en \textit{C:\textbackslash Archivos de programa} (t�picamente). Si lo est�, config�relo como se indica m�s abajo.
	\item Si no existe, vaya a la siguiente direcci�n de Internet y descargue la �ltima versi�n del  Java SE Development Kit (JDK) \cite{download-jdk}
\end{milista}

Una vez descargado, ejecute el archivo, elija el directorio de instalaci�n y siga los pasos para completar la instalaci�n. Tras instalarlo, debe configurarlo. Para ello vaya a \textit{Inicio} y haga clic derecho en \textit{Mi PC}. En las opciones, elija la pesta�a de \textit{Opciones Avanzadas} y haga clic en \textit{Variables de entorno}. En la nueva ventana que aparece, busque la variable \textit{Path} y a�ada al final \texttt{;C:\textbackslash Archivos de programa\textbackslash Java\textbackslash jdk1.x.x\_xx\textbackslash bin}, donde ''x'' hace referencia a la versi�n que tiene instalada de Java.

Ese es el directorio de instalaci�n por defecto de Java. Si en la instalaci�n eligi� otro diferente, debe introducir el que eligi� al instalarlo.

Tras esto, pruebe de nuevo a ejecutar ''java -versi�n'' en una ventana de comandos.

\subsection{Gu�a de uso}

Al ejecutar el fichero JAR correspondiente a la aplicaci�n, obtenemos una ventana como la que se muestra en la Figura \ref{fig:login}. Si la describimos de arriba abajo, se aprecian los siguientes elementos:

\begin{milista}
	\item Una caja de texto para introducir el nombre de usuario del cliente.
	\item Unos botones para elegir el rol del usuario en la aplicaci�n.
	\item Un bot�n para elegir si se desea iniciar una nueva sesi�n o unirse a una sesi�n ya iniciada. En este �ltimo caso, se debe introducir la direcci�n IP del cliente que inici� la sesi�n por primera vez, adem�s del puerto de escucha.
\end{milista}

\imagen{./imagenes/usermanual/login.png}{0.80}{Ventana de inicio de sesi�n de PlanCode}{fig:login}

Al iniciar sesi�n correctamente, se abrir� la ventana principal de PlanCode, mostradad en la Figura \ref{fig:principal}. Como se puede observar, en la lista de usuarios situada en la derecha, aparecen los usuarios conectados actualmente a una sesi�n. A continuaci�n, se muestran las distintas funcionalidades que ofrece esta herramienta:

\imagen{./imagenes/usermanual/principal.png}{0.5}{Ventana principal de PlanCode}{fig:principal}

\begin{milista}
	\item \textbf{Chat}: situado en la parte inferior de la ventana, en la primera ventana. Aqu� se mostrar�n los mensajes que los usuarios env�en, como se muestra en la Figura \ref{fig:chat}. 
	
	\imagen{./imagenes/usermanual/chat.png}{0.4}{Ventana principal de PlanCode al enviar un mensaje al chat}{fig:chat}
	
	\item \textbf{Cargar imagen}: se puede cargar una imagen en el bot�n situado en la parte superior o en la opci�n correspondiente del men� ''Archivo''. Cuando un usuario cargue una imagen, autom�ticamente se enviar� dicha imagen al resto de usuarios conectados y a los futuros clientes que se conecten, para mantener sincronizada el �rea de trabajo, como se muestra en la Figura \ref{fig:mapa}. Adem�s, como se observa en dicha figura, se actualiza el log de acciones en consecuencia.
		
	\imagen{./imagenes/usermanual/mapa.png}{0.4}{Ventana principal de PlanCode al cargar una imagen}{fig:mapa}
	
	\item \textbf{Dibujar trazo}: un cliente puede dibujar trazos libremente sobre la imagen cargada (aunque se puede dibujar tambi�n sin tener una imagen cargada) seleccionando el bot�n con el icono del l�piz, mostr�ndose dichos trazos en todos los clientes conectados. Adem�s, los futuros clientes que se conecten tambi�n recibir�n esos trazos para mantener el sincronismo de la vista compartida. En la Figura \ref{fig:trazos.png} se muestra un ejemplo de como cada cliente puede dibujar trazos, con su color correspondiente. 
	
	\imagen{./imagenes/usermanual/trazos.png}{0.35}{Ventana principal de PlanCode al dibujar trazos}{fig:trazos}
	
	\item \textbf{Eliminar trazos}: se puede eliminar un trazo, seleccionando el bot�n con el icono de la goma de borrar y pinchando en el trazo que se desea borrar. El resto de clientes se actualizar�n en consecuencia, como muestra la Figura \ref{fig:borrar}.
	
	\imagen{./imagenes/usermanual/borrar.png}{0.35}{Ventana principal de PlanCode al borrar trazos}{fig:borrar}
	
	\item \textbf{Limpiar trazos}: se pueden limpiar todos los trazos dibujados al pulsar sobre el bot�n correspondiente, o al cargar una imagen.
	
	\item \textbf{Cerrar sesi�n}: para salir del sistema, basta con darle a la opci�n correspondiente del men� ''Archivo'' o cerrar la ventana. Cuando un usuario se desconecta, se notifica en el log y se refresca la lista de usuarios conectados, como se muestra en la Figura \ref{fig:salir}
	
	\imagen{./imagenes/usermanual/salir.png}{0.5}{Ventana principal de PlanCode cuando se desconecta un usuario}{fig:salir}
	
\end{milista}




% Secci�n de ap�ndices
% Se a�ade el titulo de "Ap�ndices" al �ndice
%\addcontentsline{toc}{section}{Ap�ndices}

%\input{Apendices.tex}

% A�adimos la bibliografia al �ndice
\clearpage\phantomsection
\addcontentsline{toc}{section}{Referencias}

% Bibliograf�a. En este caso se usa BibTeX
\bibliographystyle{plain}
\pagestyle{plain} 
\bibliography{Bibliografia}

\end{document}
